%
% 35-meromorph.tex -- Meromorphe Funktionen
%
% (c) 2026 Prof Dr Andreas Müller
%
\subsection{Meromorphe Funktionen}
Funktionen mit lauter isolierten Singularitäten sind besonders
leicht zu behandeln und verdienen daher einen eigenen Namen.

\begin{definition}[meromorphe Funktion]
Eine holomorphe Funktion $f\colon U\to\mathbb{C}$,
für die $\mathbb{C}\setminus U$ aus lauter isolierten Punkten besteht,
heisst \emph{meromorph}.
\index{meromorph}%
\end{definition}

Die wichtigsten Beispiel für meromorphe Funktionen sind die rationalen
Funktionen
\[
f(z)
=
\frac{p(z)}{q(z)},
\]
wobei $p(z)$ und $q(z)$ Polynome sind.
\index{Singularitat@Singularität}%
Die Singularitäten sind genau die Nullstellen von $q(z)$ und die
Ordnung der Pole von $p(z)$ ist die Ordnung der Nullstellen von $q(z)$.
Das Residuum von $f(z)$ bei einer einfachen Nullstelle $z_0$ des Nenners
ist
\[
\operatorname{res}(z_0,f)
=
\frac{p(z_0)}{q'(z_0)}.
\]

\begin{beispiel}
Die rationale Funktion $\varphi(z) = 1/(z^2+1)$ wurde in
Beispiel~\ref{buch:analytisch:laurentreihen:bsp:cauchyverteilung}
bereits untersucht.
Man kann sie als Quotient der Polynome
$p(z)=1$ und $q(z)=z^2+1$ schreiben, die die Voraussetzungen
erfüllen: die Nullstellen $\pm i$ des Nenners sind einfach.
Es folgt, dass das Residuum
\[
\operatorname{res}(\pm i,\varphi)
=
\frac{p(\pm i)}{q'(\pm i)}
=
\frac{1}{2(\pm i)}
=
\pm\frac{1}{2i}
\]
ist, in Übereinstimmung mit dem Resultat von
Beispiel~\ref{buch:analytisch:laurentreihen:bsp:cauchyverteilung}.
\end{beispiel}

